% §1 Introduction

Simulating fermionic systems on quantum hardware requires a mapping
from the anticommuting creation and annihilation operators of second
quantization to the tensor-product structure of qubit Pauli operators.
This fermion-to-qubit encoding is a necessary middleware step between
the molecular Hamiltonian and the quantum circuit, and its choice
directly determines the Pauli weight, measurement complexity, and
circuit depth of the resulting simulation~\cite{tranter2015,
havlicek2017}.

Five encodings dominate the literature: Jordan--Wigner~\cite{jordan1928},
Bravyi--Kitaev~\cite{bravyi2002, seeley2012}, Parity~\cite{seeley2012,
bravyi2017tapering}, balanced binary tree, and balanced ternary
tree~\cite{jiang2020, miller2023bonsai}.  Each has been derived and
implemented independently, typically as a bespoke function mapping
\texttt{FermionOperator} $\to$ \texttt{QubitOperator} in software
libraries such as OpenFermion~\cite{mcclean2020}, Qiskit
Nature~\cite{qiskit2023}, and PennyLane~\cite{bergholm2022}.  Adding a
new encoding to any of these frameworks requires reimplementing the
entire mapping algorithm --- typically several hundred lines of code ---
even though the underlying algebraic structure is largely shared.

\subsection{The structural question}

This situation raises a natural mathematical question: \emph{what is the
minimal specification that uniquely determines a fermion-to-qubit
encoding?}  If such a specification exists and is compact, then the
mapping algorithm can be factored into a universal procedure
parameterized by a small data structure, and the definition of an
encoding becomes a declaration rather than a computation.

We show that two such specifications exist, each defining a class of
parameterized maps:

\begin{enumerate}
  \item \textbf{Index-set schemes.}  A triple of set-valued functions
    $\scheme = (\Uset{\cdot}, \Pset{\cdot}, \Occ{\cdot})$ from mode
    indices to subsets of qubit indices.  The universal encoding
    algorithm $\encode_\scheme$ produces Majorana operators, and thence
    ladder operators, from any valid scheme.  Jordan--Wigner,
    Bravyi--Kitaev, and Parity are each fully specified by choosing
    three such functions.

  \item \textbf{Path-based tree encodings.}  A rooted labelled tree $T$
    on $\lceil n/2 \rceil$ internal nodes, where each node has at most
    three descending links labelled $X$, $Y$, $Z$.  The encoding is
    read off from root-to-leaf paths.  Every rooted tree defines a valid
    encoding; different tree topologies recover all five known encodings
    and define infinitely many new ones.
\end{enumerate}

Both specifications factor through a common intermediate step: the
\emph{Majorana representation}, where the operators $c_j = \ad{j} +
\an{j}$ and $d_j = i(\ad{j} - \an{j})$ are expressed as single Pauli
strings.  This factorization separates the \emph{definition} of an
encoding (what makes it different from other encodings) from the
\emph{computation} of the encoding (how Majorana operators become ladder
operators become Hamiltonian terms), which is universal.

\subsection{Contributions}

This paper makes five contributions:

\begin{enumerate}
  \item \textbf{Unified framework.}  We formalize the two parameterized
    maps and prove that both are complete: every known encoding is an
    instance, and every valid parameter value yields a correct encoding
    (\cref{sec:index-set,sec:tree-encoding}).

  \item \textbf{Star-tree theorem.}  We prove that the
    Seeley--Richard--Love index-set construction satisfies the CAR if
    and only if the underlying tree is a star (depth~$\le 1$) --- a far
    sharper restriction than the monotonicity condition assumed in prior
    work.  By contrast, the path-based construction is universal.  The
    space of $n^{n-1}$ encoding trees decomposes into three
    construction-compatibility regions (\cref{sec:star-tree}).

  \item \textbf{Algebraic exactness.} We prove that the Pauli
    multiplication pipeline, when implemented with symbolic $\Zi$ phase
    tracking, preserves exactness over the Gaussian rationals $\Qi$ ---
    eliminating floating-point cancellation errors that produce spurious
    Hamiltonian terms in matrix-based approaches
    (\cref{sec:exactness}).

  \item \textbf{Tight weight bounds.}  We establish that Pauli
    weight is determined by tree depth alone:
    $w \le 2\,\depth{T} + 1$, and the balanced ternary tree achieves
    the optimal scaling $O(\log_3 n)$ (\cref{sec:tree-encoding}).

  \item \textbf{Algebraic composability.}  We show that the same
    parameterized pipeline extends to bosonic canonical commutation
    relations and mixed fermion--boson systems by abstracting the
    commutation algebra as an interface, with no change to the
    downstream Pauli multiplication (\cref{sec:extensions}).
\end{enumerate}

The framework is implemented in the open-source \textsc{FockMap}
library~\cite{fockmap2026} and validated by 497 automated tests and
eigenspectrum comparison against known molecular
benchmarks~(\cref{sec:validation}).

\subsection{Relation to prior work}

Seeley, Richard, and Love~\cite{seeley2012} introduced the update,
parity, and remainder sets for the Bravyi--Kitaev encoding and observed
that Jordan--Wigner and Parity could be expressed in the same language.
Their framework appeared to unify these three encodings under a single
index-set construction.  We show (\cref{sec:star-tree}) that this
unification is less general than it appears: the generic index-set
construction satisfies the CAR only for star trees, and the
Bravyi--Kitaev encoding actually uses a separate, Fenwick-specific
construction that is algebraically distinct.

Steudtner and Wehner~\cite{steudtner2018} identified ternary tree
structure as a unifying construction and proved $O(\log_3 n)$ weight
bounds.  Jiang et al.~\cite{jiang2020} gave the path-based
construction for optimal ternary trees, and Miller et
al.~\cite{miller2023bonsai} generalized to arbitrary tree topologies.

Our contribution builds on these foundations but shifts the emphasis
from \emph{individual encodings} to the \emph{structure of the map
itself}.  By making the parameterization explicit and the algorithm
universal, we can: (a)~define new encodings in 3--5 lines rather than
hundreds; (b)~prove properties that hold for \emph{all} encodings
simultaneously (exactness, weight bounds); (c)~extend the same
framework to non-fermionic systems; and (d)~identify structural
restrictions (such as the star-tree limitation) that are invisible when
encodings are studied one at a time.

\subsection{Outline}

\Cref{sec:prelim} establishes notation, collects core definitions,
and reviews the Majorana decomposition.  \Cref{sec:index-set} defines
the index-set parameterization and proves its correctness.
\Cref{sec:tree-encoding} defines the tree parameterization and
establishes weight bounds.  \Cref{sec:star-tree} proves the star-tree
theorem and maps the three-region structure of the encoding space.
\Cref{sec:exactness} proves algebraic exactness of the pipeline.
\Cref{sec:extensions} extends the framework to bosonic and mixed
systems.  \Cref{sec:hamiltonian} traces the complete encoding pipeline
on a worked example.  \Cref{sec:validation} presents computational
validation.  \Cref{sec:concerns} proactively addresses potential
concerns about proof technique, scope, and comparison to existing
tools.  \Cref{sec:discussion} discusses implications and open
questions.  \Cref{app:datatypes} documents the correspondence between
the algebraic data types and the mathematical structures.
